\documentclass[11pt,a4paper]{article}
\usepackage[margin=25mm]{geometry}
\usepackage[T1]{fontenc}
\usepackage{lmodern,microtype}
\usepackage{amsmath,amssymb,amsthm,mathtools}
\usepackage{booktabs,longtable,array}
\usepackage{xcolor,enumitem,fancyhdr,hyperref,xurl,needspace}
\hypersetup{colorlinks=true,linkcolor=blue!45!black,urlcolor=blue!45!black,
 pdftitle={Positive source coordinates for a smaller Diophantine consistency candidate},
 pdfauthor={Research contribution prepared for ProveIt}}
\setlength{\emergencystretch}{2em}
\setlength{\parskip}{4pt}
\setlength{\parindent}{0pt}
\clubpenalty=10000
\widowpenalty=10000
\setlist{itemsep=2pt,topsep=4pt}
\pagestyle{fancy}
\fancyhf{}
\fancyhead[L]{\small Positive source coordinates}
\fancyhead[R]{\small ProveIt research contribution}
\fancyfoot[C]{\thepage}
\setlength{\headheight}{14pt}
\newtheorem{theorem}{Theorem}
\newtheorem{lemma}[theorem]{Lemma}
\newtheorem{proposition}[theorem]{Proposition}
\theoremstyle{remark}
\newtheorem{remark}[theorem]{Remark}
\newcommand{\ZZ}{\mathbb Z}
\newcommand{\NN}{\mathbb N}
\newcommand{\PP}{\mathbb Z_{>0}}
\newcommand{\pop}{\operatorname{pop}}
\newcommand{\Con}{\operatorname{Con}}
\newcommand{\Sc}{S_{\mathrm{code}}}
\newcommand{\Tc}{T_{\mathrm{code}}}
\newcommand{\Gs}{G_{\mathrm{src}}}
\newcommand{\Us}{U_{\mathrm{src}}}
\newcommand{\Targets}{T_{\mathrm{mark}}}
\newcommand{\Gcap}{G_{\mathrm{cap}}}
\newcommand{\Kcode}{K_{\mathrm{code}}}
\newcommand{\code}[1]{\texttt{\detokenize{#1}}}
\newcommand{\basecommit}{\code{beb0192cb8dc36f78c369c080f2adfbc1887a486}}

\title{\vspace{-8mm}\textbf{Positive source coordinates for a smaller\\
Diophantine consistency candidate}\\[4pt]
\large Nine positive variables, exact degree \(1\,797\,796\),\\
465 arithmetic operations, and a compact coefficient tuple}
\author{Research contribution prepared for Vladimir Reshetnikov's ProveIt repository}
\date{23 October 2026}
\begin{document}
\maketitle
\begin{abstract}
We give an explicit refinement of the latest nine-positive-variable
construction in the reviewed ProveIt snapshot. Two positive coordinate
changes in its quadratic source reduce the weighted source from 89 to
86 records. Repacking that source lowers the largest coordinate exponent
from 116 to 100 and the convolution height from 2161 to 2047. Composing
with the existing numerical ratio and norm reduction yields one fixed
integer polynomial of exact total degree \(1\,797\,796\), evaluated by a
465-operation division-free circuit. The comparison construction has
degree \(1\,910\,404\) and 473 operations. The new equation also admits
the complete compact tuple already established in package 44, whose
four native payloads total 295,388 raw bytes.

The archive contains literal polynomial data, deterministic builders,
materialized native coefficients, exact algebra and isolation checks,
degree certificates, inherited proofs and licenses, and a secondary
15-positive-variable degree-24/89-operation refinement.
These are incremental advances. We do not establish a global simplest
equation, a new variable bound, or a construction independent of K215
and Jones. The new local equivalences and arithmetic construction are
proved here; the semantic calibration remains inherited. In particular,
this contribution does not close the repository's explicitly outstanding
object-language ZFC internalization obligation.
\end{abstract}

\section{Result, comparison, and scope}
\label{sec:scope}

The primary equation is
\begin{equation}
\boxed{\;
 F_9=\bigl[\Theta^2-\Omega^2(a_0+64a_1)\bigr]^2
       -256\Omega^4a_0a_1=0,\qquad
 (b,g,h,s,w,C,i,j,n)\in\PP^9 .
\;}
\label{eq:main}
\end{equation}
Every symbol on the right is a polynomial abbreviation defined below.
The only unknowns are the nine displayed inputs. All exponents are
fixed nonnegative integers. No variable exponent, divisibility
condition, bitwise operation, root, or quantifier is hidden in the
equation. Such operations occur only in its proof.

The primary machine-readable specification is
\path{research_nine/compact_candidate/generated/gap_100_positive.json}.
Every nonliteral fixed leaf has a polynomial recipe in four specified
integers, or is a specified power of two. The native integers are
supplied as authenticated unsigned little-endian payloads and by a
deterministic reconstruction program. Thus no universal-polynomial
parameter remains unspecified.

The baseline is package 59's \code{nine_free_adjunct} construction at
the repository snapshot \basecommit. We use its cost convention:
one operation for every binary addition, subtraction or multiplication,
including scalar multiplications and the multiplications implementing
fixed exponents of polynomial registers. Reuse and access to fixed
integer coefficients are free. This is
arithmetic circuit size, not bit complexity.

\begin{table}[htbp]
\centering
\small
\begin{tabular}{@{}lrrrr@{}}
\toprule
Construction & Positive variables & Exact degree & Operations & Port bytes\\
\midrule
Package 59, free carrier & 9 & 1,910,404 & 473 & 533,099\\
New source, original tuple & 9 & 1,797,796 & 465 & 533,099\\
New source, compact tuple & 9 & 1,797,796 & 465 & 295,388\\
Package 58, integral-root form & 15 & 24 & 90 & 295,388\\
Supplement: shifted transport & 15 & 24 & 89 & 295,388\\
\bottomrule
\end{tabular}
\caption{Measured improvements. Raw bytes count the four native
ports, not the expanded polynomial's coefficients. The compact
native tuple is inherited from package 44. Its use here is a new
coherent composition, not a new native coefficient-size record.}
\label{tab:metrics}
\end{table}

The degree decreases by \(112\,608\), approximately \(5.89446\%\);
the circuit saves eight operations, approximately \(1.69133\%\).
Using the compact tuple reduces native raw bytes by approximately
\(44.5904\%\). Both complete specializations are retained separately.
The natural-domain translation
\[
 F_9^{\NN}(x_1,\ldots,x_9)=F_9(x_1+1,\ldots,x_9+1),
 \qquad \NN=\{0,1,\ldots\},
\]
has the same exact degree and a 472-operation supplied circuit.
No nine-unrestricted-integer claim is made.

\subsection{The proof boundary}

There are three different levels.
\begin{enumerate}
\item \textbf{New algebra and encoding.}
The positive source substitutions, the 86-record identities, finite
code, range and valuation calibration, and exact circuit metrics are
established here. With the explicitly stated numerical Jones ratio
theorem, they identify positive zeros of \(F_9\) with common positive
zeros of the specified quadratic source.
\item \textbf{Inherited semantics.}
The native compiler, Pell arguments, Wang943 language and calibrated
K215 endpoint identify the intended source with a finite syntactic
ZFC refutation. Those are inherited mathematical dependencies. We
checked the new adapter and literal native data, not the entire
transitive semantic proof chain.
\item \textbf{Internal ZFC derivation.}
The reviewed repository's \code{OBJECT_ZFC.md} explicitly leaves
the actual internal equivalence derivation open. An external Lean
theorem about the \code{BProv} predicate is not itself a derivation
of that theorem inside the object calculus. We supply neither a new
object-language derivation nor a new end-to-end Lean/Rocq certificate.
\end{enumerate}

The fully established new result is therefore a smaller explicit
arithmetic candidate preserving the inherited mathematical semantics.
If the inherited source calibration and the imported arithmetic/compiler
correctness statements are established internally in ZFC, the
reductions transport them to the intended statement
\begin{equation}
\mathrm{ZFC}\vdash
 \left[
 \neg\exists\mathbf x\in\PP^9\;F_9(\mathbf x)=0
 \ \longleftrightarrow\ \Con(\mathrm{ZFC})
 \right].
\label{eq:internal-target}
\end{equation}
This report does not present \eqref{eq:internal-target} as a newly
exported internal theorem. A missing machine certificate alone would
not refute an ordinary mathematical in-ZFC proof; an external
standard-natural equivalence alone does not supply one either.
No consistency assumption is used, and neither side of the equivalence
is asserted to be provable in ZFC.

\section{The fixed compact coefficient instance}
\label{sec:coefficients}

Use the entire native tuple from package 44:
\[
 \beta=B_\ast-1,\qquad K_0=D_C+B_\ast D_R,\qquad
 M_C,\qquad M_F=M_F^{\mathrm{native}}+B_\ast-1,\qquad
 B_\ast=2^{590069}.
\]
The native radix \(B_\ast\) is a fixed integer, distinct from the
packing radix \(B=a_\ast b^2\) introduced later. The full native
proof and data are in \path{upstream44/native/}.

\begin{center}
\begin{tabular}{@{}lrrr@{}}
\toprule
Port & Bit length & Raw bytes & Gzip bytes\\
\midrule
\(\beta\) & 590,069 & 73,759 & 108\\
\(K_0\) & 592,884 & 74,111 & 8,940\\
\(M_C\) & 590,069 & 73,759 & 193\\
\(M_F\), already shifted & 590,070 & 73,759 & 204\\
\midrule
Total & & 295,388 & 9,445\\
\bottomrule
\end{tabular}
\end{center}
The binding aliases \code{K} to \(K_0\), and \code{MF_source} to
\(M_F\). No individual port may be substituted from another tuple.

For an explicit finite recipe, put \(\rho=2^{67}\), \(\mu=2^{51}-2\).
Enumerate the literal ordered 943-tile list by \(i=0,\ldots,942\).
For tile \(i=(w_i,n_i,e_i,s_i)\), set
\[
 v_i=w_i+2^8n_i+2^{17}e_i+2^{25}s_i+127\cdot2^{34},
 \qquad c_i=1+2^{10}(2^{41}-v_i).
\]
Adjoin copy positions \(1023+j,1013+j,1003+j,983+j\), with respective
bit-position lists
\[
\begin{split}
 &(0,\ldots,7,34),\quad(8,\ldots,16,35),\\
 &(17,\ldots,24,36),\quad(25,\ldots,33,37).
\end{split}
\]
If a copy position has bit position \(p\), set its \(c_e=2^{10+p}\).
Adjoin \(c_{1095}=2^{48}\), \(c_{2195}=2^{49}\),
\(c_{3296}=2^{50}\), and \(c_{1063}=c_{1064}=0\).
Let \(\mathcal E\) be the resulting 986-position support, including
the two zero-valued positions. Then
\begin{align*}
 D_C={}&\rho^{30}+\rho^{62}+\rho^{1100}+\rho^{2201}
       +\sum_{e\in\mathcal E}c_e
        (\rho^{4393-e}+\rho^{5495-e}),\\
 D_R={}&\rho^{42},\\
 Z_\ast={}&(2^{65}-1)\sum_{e=8792}^{8804}\rho^e
                         +(2^{19}-1)\rho^{8805},\\
 M_C={}&B_\ast-1-\sum_{\substack{e\in\mathcal E\\e\ne1}}\rho^e
                         -2\rho^{1063},\\
 M_F^{\mathrm{native}}={}&4+\mu(\rho^{4393}+\rho^{5495})
       +\sum_{j=0}^{9}\rho^{1013+j}
       +\sum_{j=0}^{8}\rho^{1065+j}
       +\rho^{2195}+\rho^{3296}+Z_\ast .
\end{align*}
The tile list is fixed archive data, not an existentially chosen program.
The native builder reconstructs it from the 430 literal K215 transitions.

Replaying the independent native auditor checks 12 pinned source files,
66,062 native basis interactions, the literal transitions,
38,663 one-bit mutations, the rank controller and exact payloads.
Its 87,380 small complement models supplement the unbounded proof.
The rank-control prime is \(581197\), with dimensions
\(67\cdot8807=590069\). These finite checks do not replace the
native semantic theorem.

\section{Normalize positive source coordinates before packing}
\label{sec:source}

The source has 31 positive coordinates and 22 quadratic residuals.
Names in this section refer to the source, independently of similar
names in the final numerical tail. The new coordinate order is
\begin{equation}
\begin{split}
 &(J,F,\alpha,z,f,h,i,o,j,s,w,t,\eta,\zeta,y,\gamma,\\
 &\hspace{13mm}N,E,p,c,a,X,D,C_2,S,F_2,V,g,I,B_{\mathrm{aux}},\Delta).
\end{split}
\label{eq:source-order}
\end{equation}
Here \(t,p,S\) replace \(\tau,P,L\), respectively.

\subsection{Two Pell coordinates become positive gaps}

Let \(k=\eta+\zeta\). The old source includes
\[
 a-E=sN,\quad c-\eta=k(a-E),\quad
 P-k=E(c-\eta),\quad \tau^2=P(P-k)+1.
\]
The first two rows force \(a-E>0\), \(c-\eta>0\). Consequently
\[
 p:=P-k=E(c-\eta)>0,\qquad
 \tau^2=p^2+kp+1>p^2,\qquad t:=\tau-p>0.
\]
The changed rows are
\begin{equation}
\boxed{\quad p-Ec+\eta E=0,\qquad
 t^2+2pt-\eta p-\zeta p-1=0.\quad}
\label{eq:gap}
\end{equation}
Conversely, positive \(p,t\) give positive
\(P=p+\eta+\zeta\), \(\tau=p+t\), satisfying precisely the old rows.
Neither old variable occurs elsewhere. This is a bijection of the
complete positive common zero sets, not a restriction to selected
canonical witnesses.

The first row loses two weighted monomial records. The second retains
five, but replaces the active square \(P^2\) by the pair \(pt\).
That change in the active pair pattern enables a better packing.
Coordinate normalization before packing is the main new ingredient.

\subsection{Supply the square instead of the square minus one}

Write \(q=\beta J+1\). Replace the old register \(L=q^2-1\)
by \(S=L+1\). The first two rows become
\begin{equation}
 S-\beta^2J^2-2\beta J-1=0,\qquad
 N-S-\beta JS=0.
\label{eq:square-shift}
\end{equation}
The inverse is positive: the first row gives \(S=q^2\ge4\),
so \(L=S-1>0\).

Retain the inherited invertible row operations
\[
 \rho_{14}=r_{14}+r_2,\qquad
 \rho_{15}=r_{15}+r_3,\qquad \rho_j=r_j\ (j\ne14,15).
\]
Under \(L=S-1\), row 1 gains one record, row 2 loses two, and
row 14 keeps its count. Its \(-FL\) contributes \(+F\), cancelling
the old \(-F\), while \(-\alpha L\) contributes \(+\alpha\).
No quadratic record changes in this second step. Thus
\[
 89\longrightarrow87\longrightarrow86
 =7\text{ constant}+36\text{ linear}+43\text{ quadratic records},
\]
representing 39 distinct quadratic monomials.

\subsection{The transport convention and coherent rebinding}

The current package-31 source and native package 44 use the shifted
positive transport coordinate \(z_+\):
\[
 (K_0+w)C=F+z_+(q-1),\qquad C=q-F-\alpha.
\]
Our weighted source uses \(z=z_++1\):
\[
 (K_0+w)C+(q-F)-z(q-1)=1.
\]
The positive inverse uses inherited source consequences
\(C>0\), \(F<q\), and \(w>q\), giving
\[
 (z-1)(q-1)=(K_0+w)C-F>0,\qquad z\ge2.
\]
This convention is explicit in the independent checker: it substitutes
\(z_+=z-1\), accounting for \(+\beta J\) in row 15. No historical
byte identity between these two source presentations is claimed.

Package 44, Section 7, proves the same shifted 16-variable Pell system
for the compact tuple. Projection of the 31-variable source forgets
defining registers. Its canonical Pell witnesses give the reverse
positive lift. In particular the delicate register
\(g=y^2-V^2\) obeys
\[
 (I^2-1)g=V^2-1,
\]
and the canonical witnesses have \(I,V>1\), so \(g>0\).
The complete native theorem supplies this lift; it is not assumed
for arbitrary assignments without the canonical completeness argument.

\section{Literal degree-separated packing}
\label{sec:packing}

The literal 86 weighted records are in
\path{research_nine/compact_candidate/data/parent_source.json}.
All 22 residual formulas, their targets and coordinate weights are
also displayed in Appendix~\ref{app:source}.
Their independent reconstruction treats all four ports as independent
symbols. In the coordinate order \eqref{eq:source-order}, the weights
are
\begin{equation}
\begin{split}
 (w_1,\ldots,w_{31})={}&
 (6,10,37,46,4,83,60,50,65,45,20,87,96,80,0,41,\\
 &3,89,95,100,97,84,1,67,21,59,18,19,61,14,64).
\end{split}
\label{eq:weights}
\end{equation}
The zero weight is intentional. Assign target exponents by the paired
lists
\begin{align*}
 \text{rows: }&
 (2,18,9,20,19,7,12,17,11,10,1,15,13,21,4,16,8,3,22,5,14,6),\\
 \text{targets: }&
 (27,72,270,320,448,561,616,686,887,893,906,966,\\[-2pt]
 &\hspace{13mm}1139,1190,1255,1279,1472,1499,1545,1733,1822,2016).
\end{align*}
Let \(q_j\) be the target of row \(j\), and \(W=100\).
The targets are unrelated to the source abbreviation \(q=\beta J+1\).

If \(\rho_j(\mathbf y)=\sum_\alpha a_{j,\alpha}\mathbf y^\alpha\),
set
\[
 c_{j,\alpha}=
 \begin{cases}
 2a_{j,\alpha},&|\alpha|=0,1,\\
 a_{j,\alpha}\prod_i\alpha_i!,&|\alpha|=2,
 \end{cases}
\]
and define the three kernels
\begin{equation}
 D_d(T)=\sum_{j,\ |\alpha|=d}
  c_{j,\alpha}T^{q_j-\sum_i\alpha_iw_i},\qquad d=0,1,2.
\label{eq:kernels}
\end{equation}
All exponents are nonnegative. Expanded lists are supplied in
\code{data/split_kernels.json}.

\begin{lemma}[Exact target extraction]
\label{lem:extract}
For \(U_{\mathbf y}(T)=\sum_i y_iT^{w_i}\),
\[
 [T^{q_j}]\bigl(D_2(T)U_{\mathbf y}(T)^2+
               D_1(T)U_{\mathbf y}(T)+D_0(T)\bigr)=2\rho_j(\mathbf y).
\]
The kernel degrees are \((1545,1920,1847)\), the leading coefficient
of \(D_2\) is \(1\), and
\(\Kcode=\max_d(\deg D_d+dW)=2047\).
\end{lemma}
\begin{proof}
For each source record and every monomial of its degree, the checker
tests whether the shifted exponent reaches any target. The degree
\(0,1,2\) bases have \(1,31,496\) monomials. All
\(7+36\cdot31+43\cdot496=22451\) interactions are enumerated.
A hit occurs only for the record's own row and monomial.
The factorial normalization compensates for the two copies of cross
terms in a square. Exact coefficient-polynomial comparison proves
the identities, and the literal nonzero kernel heads give the degrees.
\end{proof}

\subsection{Packing registers and bounds}

Every weighted coefficient satisfies \(|c_{j,\alpha}|<M\), with
\[
 M=2^{1182955},\quad a_\ast=4096M=2^{1182967},\quad
 c_\ast=2^{9463725}.
\]
The exponents come from every exact coefficient recipe and the four
port bit bounds. For positive outer \(b,g\), define
\begin{align}
 B&=a_\ast b^2,& \ell&=\sum_iB^{w_i},&\Us&=\ell+g,\nonumber\\
 \Gs&=(D_2(B)\Us+D_1(B))\Us+D_0(B),&
 \Targets&=\sum_jB^{q_j},&N_0&=B^{101},\nonumber\\
 \mathrm{Mask}&=N_0-1-(b-1)\ell,&
 J_{\mathrm{code}}&=2\Gs+\Targets,\nonumber\\
 \Sc&=g+N_0J_{\mathrm{code}},&
 \Tc&=\mathrm{Mask}+N_0(B-2)\Targets .\label{eq:pack-registers}
\end{align}
Finally set
\begin{equation}
\begin{aligned}
 \mathcal L&=bB^{2149},&
 R&=\mathcal L(\Sc+\Tc+1)-\Sc-1,\\
 N&=c_\ast bB^{1128},&
 t_{\mathrm{cap}}&=bB^{100},\qquad
 \Gcap=32t_{\mathrm{cap}} .
\end{aligned}
\label{eq:calibration}
\end{equation}

\begin{lemma}[Unconditional positivity and the full soundness range]
\label{lem:bounds}
For every positive \(b,g\), \(\Gs>0\), \(\Sc,\Tc\ge0\),
\(R,N\ge8\), and \(b\le R,N\). If \(g<4bB^{100}\), then
\[
 0<g<N_0,\qquad \Sc+\Tc<B^{2149}\le\mathcal L.
\]
\end{lemma}
\begin{proof}
Distinct kernel exponents give
\[
 D_2(B)>\tfrac12B^{1847},\quad
 |D_1(B)|<2MB^{1920},\quad |D_0(B)|<2MB^{1545}.
\]
Since \(\Us\ge B^{100}\), the negative contributions relative to
\(B^{1847}\Us^2\) are below \(2M/B^{27}\) and \(2M/B^{502}\),
each below \(1/8\). Hence \(\Gs>\tfrac14B^{1847}\Us^2>0\).
The mask has digits \(B-b\) at selected weights and \(B-1\)
elsewhere, so lies in \([0,N_0)\).
Also \(R=(\mathcal L-1)\Sc+\mathcal L\Tc+\mathcal L-1\ge\mathcal L-1\).

On the bounded branch, \(\ell<2B^{100}\) and \(\Us<6bB^{100}\).
With \(Z=B^{2048}\), geometric summation gives
\[
 |\Gs|<\frac{43Mb^2Z}{B-1},\qquad
 \frac{2\Gs}{Z}<\frac{86Mb^2}{B-1}<\frac1{32}.
\]
The maximum target is 2016, so
\(\Targets<Z/64\), \((B-2)\Targets<Z/64\).
As \(4b<B\), \(g<N_0\), and
\(J_{\mathrm{code}}+(B-2)\Targets+2<Z\).
It follows that \(\Sc+\Tc<N_0Z=B^{2149}\).
This proves the whole range delivered by soundness, rather than
only the narrower completeness range.
\end{proof}

\section{From digit disjointness to one binomial valuation}
\label{sec:binary}

Assume now \(b=2^k\), allowing \(k=0\), and \(B=2^v\), where
\(v=1182967+2k\). On the bounded branch, the low-mask condition
\(g\mathbin{\&}\mathrm{Mask}=0\) is equivalent to
\[
 g=\sum_i z_iB^{w_i},\qquad 0\le z_i<b.
\]
Put \(y_i=z_i+1\). Then \(1\le y_i\le b\) and
\(\Us=\sum_i y_iB^{w_i}\).
Every coefficient \(h_u\) of the source convolution satisfies
\[
 |h_u|<M\bigl((\sum_i y_i)^2+\sum_i y_i+1\bigr)
       \le993Mb^2<B/4.
\]
At a target, the signed base-\(B\) digit of \(J_{\mathrm{code}}\)
is \(2h_u+1\), with incoming borrow \(-1\) or \(0\).
Its normalized digit avoids the mask \(B-2\) precisely when
\(h_u=0\): zero gives digit 0 or 1, positive \(h_u\) gives a digit
between 2 and \(B/2\), and negative \(h_u\) gives a digit above \(B/2\).
Lemma~\ref{lem:bounds} ensures finite, nonoverlapping words.
Lemma~\ref{lem:extract} therefore proves
\begin{equation}
 \Sc\mathbin{\&}\Tc=0
 \quad\Longleftrightarrow\quad
 g\text{ encodes a positive common source zero}.
\label{eq:decode}
\end{equation}
Conversely, every positive source zero is encodable for every
sufficiently large binary \(b\), with \(0<g<bB^{100}\), hence
\(32g<\Gcap\). The strict lower bound follows because the all-one
source tuple is not a zero: its fifth residual is 1.

For \(\mathcal L=2^u\) and \(\Sc+\Tc<\mathcal L\), the two blocks
\[
 R=(\Sc+\Tc)\mathcal L+(\mathcal L-1-\Sc)
\]
give
\[
 \pop(R)=u+\pop(\Tc)-\tau_{\mathrm{carry}},\qquad
 \tau_{\mathrm{carry}}
 =\pop(\Sc)+\pop(\Tc)-\pop(\Sc+\Tc)\ge0.
\]
The defect vanishes exactly when \(\Sc\mathbin{\&}\Tc=0\).
Counting the low mask and the 22 target masks gives
\[
 \pop(\Tc)=123v-31k-22,\qquad u=k+2149v.
\]
Consequently
\begin{align*}
 u+\pop(\Tc)&=2272v-30k-22\\
 &=2(9463725+k+1128v)=2\log_2N,
\end{align*}
using \(2\cdot9463725=16\cdot1182967-22\).
Legendre's formula
\(\nu_2\binom{2R}{R}=\pop(R)\) then yields
\begin{equation}
 N^2\mid\binom{2R}{R}
 \quad\Longleftrightarrow\quad\Sc\mathbin{\&}\Tc=0.
\label{eq:binomial}
\end{equation}
The boundary \(b=1\) is included, but cannot produce a final zero:
then \(\mathrm{Mask}=N_0-1\), while \(0<g<N_0\).
The calibration makes no unproved restriction to a sparse family
of binary bases.

\section{The complete nine-variable polynomial}
\label{sec:tail}

For the nine positive inputs in \eqref{eq:main}, put \(j_0=j-1\ge0\).
Use \eqref{eq:pack-registers}--\eqref{eq:calibration}, and define
\begin{align}
 U&=N^2w,&Y&=N^2s,&M_0&=RY,\nonumber\\
 P&=2M_0^2U,&K&=R+1+h(P-1),&A&=M_0(U+1),\nonumber\\
 B_1&=2R+1,&W_1&=bw,&J&=4A-5,\nonumber\\
 Q&=3W_1C-2(W_1^2-1),&
 D&=(A^2-1)C^2+1,&k_0&=iJC,\nonumber\\
 E&=k_0DC,&F&=(A^2-1)E^2+1,\nonumber\\
 G&=A+F(CD+1-A),&
 H&=C+F(B_1+j_0C),\nonumber\\
 I&=(G^2-1)H^2+1,&
 a_0&=(P^2-1)K^2+1,&a_1&=DFI.
\label{eq:tail-registers}
\end{align}
Set
\begin{align}
 m&=K^2(4t_{\mathrm{cap}}-g)
       -16t_{\mathrm{cap}}(C-KY)^2,\nonumber\\
 V_0&=PK,&T_{\mathrm{pair}}&=k_0(8D^2-4D-1),&
 V_1&=T_{\mathrm{pair}}GH,\nonumber\\
 \Omega&=Jm,&E_0&=J(n+W_1^2)+Q,&
 \Theta&=E_0-\Omega(V_0+V_1).
\label{eq:norm-registers}
\end{align}
Then \(F_9\) is exactly \eqref{eq:main}. The fixed kernels, source
records, coefficient tuple and the displayed definitions specify
the entire polynomial. No intermediate register adds a witness.

\subsection{The inherited numerical ratio interface}

We use the numerical Jones ratio soundness and completeness theorem
in the exact reparameterized form retained by package 59.
Under the unconditional hypotheses in Lemma~\ref{lem:bounds}
and the recovered carrier range \(C>B_1+W_1\), the conditions
\[
 a_0\text{ square},\qquad a_1\text{ square},\qquad J\mid Q,\qquad m>0
\]
imply binary \(b\), \(8g<\Gcap\), and
\(N^2\mid\binom{2R}{R}\). Conversely, the encoded source witnesses
with \(32g<\Gcap\) admit the required ratio witnesses.
This is an imported unbounded arithmetic theorem, not a consequence
of the finite audits.

Its CRT coordinate substitution is exact. The previous condition is
\[
 JF\mid QF+J(H_{\mathrm{old}}-C),\qquad
 H_{\mathrm{old}}=B_1+j_{\mathrm{old}}C,\quad j_{\mathrm{old}}\ge1 .
\]
Since \(F\equiv1\pmod J,\pmod C\), it is equivalent to \(J\mid Q\)
and \(F\mid H_{\mathrm{old}}-C\).
The positive quotient is \(B_1+j_0C\) for a unique \(j_0\ge0\).
Conversely,
\[
 j_{\mathrm{old}}=1+\frac{F-1}{C}B_1+Fj_0
\]
is a positive integer giving exactly the displayed \(H\).
The division occurs only in the proof; \(C\mid F-1\).
The positive input must therefore be \(j=j_0+1\), rather than
silently imposing \(j_0>0\). The normalized margin satisfies
\[
 8m=[K^2-4(C-KY)^2]\Gcap-8gK^2.
\]

\section{The norm has no extraneous positive zeros}
\label{sec:norm-proof}

The logical order is crucial: establish the sign of \(m\) before
assuming a carrier range or integral roots.

\begin{lemma}[Unconditional strict majorants]
For every positive input assignment,
\[
 V_0^2-a_0=K^2-1>0,\qquad
 T_{\mathrm{pair}}^2-64DF=(8D+1)k_0^2-64D>0,
\]
and
\begin{equation}
\begin{split}
 \Delta_1:=V_1^2-64a_1
 ={}&[(8D+1)k_0^2-64D](GH)^2\\
   &+64DF(H^2-1)>0.
\end{split}
\label{eq:defect}
\end{equation}
Also \(E_0>0\).
\end{lemma}
\begin{proof}
The unconditional bounds imply
\(A\ge2\), \(J\ge3\), \(D\ge A^2\), \(D>C^2\),
\(F\ge A^2\), \(G>A\), \(H>1\), and \(k_0\ge3\).
For the second identity substitute
\(F=1+k_0^2D^2(D-1)\) and expand
\((8D^2-4D-1)^2\). Its difference is at least \(8D+9\).
The third identity follows from \(I=(GH)^2-H^2+1\).
Finally \(A>W_1\), and
\[
 E_0=Jn+(J-2)W_1^2+3W_1C+2>0,
\]
including assignments where the free carrier makes \(Q\) negative.
\end{proof}

\begin{theorem}[Arithmetic zero equivalence]
\label{thm:arith}
With the coherent fixed coefficients and the stated numerical ratio
theorem, \(F_9\) has a positive zero exactly when the new
31-positive-coordinate quadratic source has a common zero.
\end{theorem}
\begin{proof}
Suppose \(F_9=0\). If \(m=0\), then \(\Omega=0\) and
\(F_9=E_0^4>0\), a contradiction. Otherwise the real norm
factorization supplies signed real roots
\(r_0^2=a_0\), \(r_1^2=64a_1\) with
\[
 n+W_1^2+\frac QJ=m(V_0+V_1+r_0+r_1).
\]
The left side is positive. The strict majorants make the right
parenthesis positive for every sign choice. Thus \(m>0\), without
an integrality assumption.

The margin now gives \(g<4t_{\mathrm{cap}}\) and
\(|C-KY|<K/2\). As \(Y\ge64\),
\[
 C>K(Y-\tfrac12)>2K\ge2(P+R)>2R+1+W_1.
\]
Here \(K\ge P+R\) and \(P>U\ge W_1\).
Also \(A^2>P\), \(D>C^2\), \(F\ge A^2\), \(G\ge A\),
and \(I\ge G^2\), giving
\[
 a_1=DFI>C^2A^4>K^2P^2>a_0 .
\]
The rational sum \(x=r_0+r_1=\Theta/\Omega\) is nonzero.
The identity \(2xr_0=x^2+a_0-64a_1\) makes both roots rational,
hence integral because their squares are integers.
Moreover \(r_1/8\) is a rational square root of \(a_1\), hence an
integer. The factor equation implies \(J\mid Q\). Ratio soundness
now supplies binary \(b\), \(8g<\Gcap\), and binomial divisibility.
Equations \eqref{eq:binomial} and \eqref{eq:decode} recover the
positive source zero.

Conversely, encode a positive source zero with sufficiently large
binary \(b\), giving \(32g<\Gcap\) and the binomial condition.
Ratio completeness supplies witnesses and integral roots.
Transport the CRT coordinate as above and choose positive roots.
Since \(W_1^2<D\), \(W_1C<D\), \(Q/J\le W_1C\), and
\(T_{\mathrm{pair}}\ge9D\), one has
\[
 W_1^2+\frac QJ<2D<V_1 .
\]
Thus
\[
 n=m(V_0+V_1+r_0+r_1)-W_1^2-\frac QJ
\]
is a strictly positive integer. It makes one norm factor vanish,
so \(F_9=0\) in the stipulated domain.
\end{proof}
The source changes of Section~\ref{sec:source} and the coherent native
adapter transport this equivalence to the inherited semantics, at
the proof level stated in Section~\ref{sec:scope}.

\section{Exact degree, rather than naive degree propagation}
\label{sec:degree}

The main circuit has material leading cancellation. The certificate
uses the exact identity
\begin{equation}
\begin{split}
 \Theta^2-\Omega^2(a_0+64a_1)
 ={}&E_0^2-2E_0\Omega(V_0+V_1)\\
  &+\Omega^2[(K^2-1)+\Delta_1+2V_0V_1].
\end{split}
\label{eq:cancellation}
\end{equation}
Fixed integers have degree zero. On the all-inputs diagonal, the
key source degrees and heads are:
\begin{center}
\begin{tabular}{@{}lrl@{}}
\toprule
Register & Exact degree & Diagonal leading coefficient\\
\midrule
\(\Us\) & 200 & \(a_\ast^{100}\)\\
\(\Gs\) & 4094 & \(a_\ast^{2047}\)\\
\(\Sc\) & 4296 & \(2a_\ast^{2148}\)\\
\(\Tc\) & 4236 & \(a_\ast^{2118}\)\\
\(\mathcal L\) & 4299 & \(a_\ast^{2149}\)\\
\(R\) & 8595 & \(2a_\ast^{4297}\)\\
\(N\) & 2257 & \(c_\ast a_\ast^{1128}\)\\
\bottomrule
\end{tabular}
\end{center}
The decisive tail degrees are
\[
\begin{array}{c|rrrrrrr}
 &A&D&k_0&F&G&H&a_0\\ \hline
 \deg&17625&35252&17627&141010&176263&149605&122942
\end{array}
\]
and
\[
\begin{array}{c|rrrrrr}
 &a_1&m&\Omega&V_0&V_1&\Delta_1\\ \hline
 \deg&827998&70703&88328&61471&413999&722242.
\end{array}
\]
The term \(\Omega^2\Delta_1\) uniquely dominates
\eqref{eq:cancellation}, at degree 898898. Its square has degree
1797796, whereas the final subtracted term
\(256\Omega^4a_0a_1\) has degree only 1304252.
The exact final diagonal leading coefficient is
\begin{equation}
 \boxed{2^{186}a_\ast^{898488}c_\ast^{400}>0.}
\label{eq:head}
\end{equation}
The certificate supplies a matching multivariate upper bound, and
the nonzero diagonal term attains it. The total degree is therefore
exactly \(1\,797\,796\).

For this calibrated family,
\[
 \deg F_9=104\deg R+400\deg N+8W+316
        =816\Kcode+1224W+5044 .
\]
The reduction decomposes as
\[
 816(2161-2047)+1224(116-100)=112608 .
\]
This gives a useful search objective: coordinate design and packing
height should be optimized together, since reducing the monomial
count alone need not reduce the final degree.

The primary positive evaluator has 159 additions, 22 subtractions
and 284 multiplications, all live. The cancellation-resolving
circuit is a separate certificate, not the counted primary evaluator.
The natural translation has 167 additions, 21 subtractions and
284 multiplications. Its internal \(j_0\) is just the natural
\(j\)-input, accounting for the subtraction saved.

\section{A small degree-24 companion}
\label{sec:low}

The companion has 15 positive unknowns, exact degree 24 and
89 operations, compared with package 58's 90.
The substantial degree reduction is inherited from package 58;
this contribution saves one transport operation.
The old shared-register residual is
\[
 T_{\mathrm{old}}
 =q\{K_0C+(q-F)-z_{\mathrm{old}}(q-1)\}+XC-q.
\]
Putting \(z_{\mathrm{old}}=z+1\) gives exactly
\begin{equation}
 \boxed{T_{\mathrm{new}}=C(qK_0+X)-q\{F+z(q-1)\}.}
\label{eq:transport-low}
\end{equation}
The inherited source inequalities give the positive inverse via
\((z_{\mathrm{old}}-1)(q-1)=(K_0+w)C-F>0\).
An eight-operation residual becomes a seven-operation residual.
The whole sum of squares has 47 multiplications, 27 additions and
15 subtractions. Its unchanged degree-12 residual ensures exact
degree 24 after squaring. The supplement includes both literal
circuits, a complete proof and an exact symbolic identity/degree audit.
This is not a degree-22 result or a variable-count improvement.

\section{Alternative substrates and the limits of the present search}
\label{sec:alternative-substrates}

The construction problem does not logically require a particular Turing
machine, nor a particular universal-polynomial presentation.  We examined
alternatives to the inherited substrate as well as algebraic changes to its
final arithmetic kernel.  None of the independent-substrate routes examined
here produced a smaller, fully instantiated polynomial with a complete
correctness argument.  This is a statement about the work completed in this
iteration, not an impossibility theorem or a claim that the existing
substrate is optimal.

\subsection{A direct proof-checker route}

Fix an effective proof calculus for ZFC and an explicit coding of its
formulas and finite proofs.  The predicate
\(\operatorname{Ref}_{\mathrm{ZFC}}(p)\), asserting that \(p\) codes a
refutation, can be chosen primitive recursive.  A direct compiler could
therefore replace the machine simulation by a polynomial
\(D(p,\mathbf x)\) with a verified representation theorem
\[
 \mathrm{ZFC}\vdash
 \bigl(\exists\mathbf x\in\mathbb N^m\;
       D(p,\mathbf x)=0\bigr)
 \longleftrightarrow \operatorname{Ref}_{\mathrm{ZFC}}(p).
\]
After existentially quantifying \(p\), the absence of zeros would be
equivalent to syntactic \(\operatorname{Con}(\mathrm{ZFC})\).  This is a
valid architecture that makes no reference to K215.

To turn that architecture into the requested fixed equation, the following
data and arguments must actually be supplied.  The syntax encoding must
specify formula parsing, bound-variable handling and substitution; axiom
recognition must include the Separation and Replacement schemes; and the
proof checker must specify every allowed inference.  The arithmetic
compiler must then emit polynomial syntax, literal fixed coefficients and
a proof relating its integer witnesses to the checker.  Finally, the
checker must be calibrated against the same effective ZFC calculus used in
the consistency assertion.  It is the combined cost of these components
that should be compared with the present construction.

The MRDP theorem establishes the availability of a Diophantine
representation, but an unspecified representation or an unspecified
universal-polynomial parameter is not an explicit coefficient instance.
Likewise, a theorem accepting supplied representation syntax does not by
itself construct that syntax for a particular proof checker.  The current
ProveIt \texttt{QUARTIC\_COMPILER.md} explicitly distinguishes these two
interfaces.  Closely related distinctions between mathematical
construction, polynomial data and formal verification also occur in the
Isabelle universal-pair development.\footnote{Jonas Bayer and Marco David,
\emph{A Formal Proof of Complexity Bounds on Diophantine Equations},
\url{https://arxiv.org/abs/2505.16963}; the associated development is
\url{https://isa-afp.org/entries/Diophantine_Universal_Pairs.html}.}

A finite logical presentation or a smaller equiconsistent theory could be
a useful source language.  However, counting axioms does not measure the
size of the compiled polynomial.  Such a proposal must provide effective
proof translations sufficient to justify the intended consistency
equivalence, and must still pay for its syntax checker.  Replacing proofs
by finite models of ZFC would not work: ZFC itself proves that its Infinity
axiom rules out finite models.  Replacing the theory by an assertion that
a model of ZFC exists also requires a separate strength analysis; that
assertion cannot simply be treated as an inexpensive equivalent axiom.

\subsection{Generic degree reduction is not a new small quartic}

An arithmetic circuit can be converted to quadratic equations by
introducing values for intermediate registers, and a sum of their squares
then has degree at most four.  This exchanges degree for unknowns.  Our
straightforward lift of the ordinary Pell kernel did not improve the
existing source-31 point: that source already has 31 positive unknowns and
22 quadratic residuals.  Its comparison evaluator has 60 operations and
its quartic sum of squares has 125 operations.  These figures are the
existing \texttt{31-sun10/quartic31/QUARTIC60\_PROOF.md} benchmark, not a
new result of this report.

One illustrative failed elimination concerns \(X=wq\).  Combining it
with the transport relation can remove the letter \(w\) from a polynomial
identity, but generally leaves only \(q\mid XC\).  Restoring the old
witness \(w=X/q\) would require \(q\mid X\).  It is invalid to cancel
\(C\) without first establishing \(\gcd(q,C)=1\).  In the contemplated
route, the desired coprimality depended on source-decoding consequences
that themselves used \(q\mid X\).  We therefore did not accept this as a
reduction in the number of variables.  This identifies a missing argument
in that particular route; it does not rule out a different elimination.

\subsection{Two auxiliary congruence routes that do not lower degree}

The degree-24 integral-root construction has an auxiliary Pell norm
\[
 U V^2-(U-1)y^2=1,
 \qquad U=I^2=\Delta(f^2-1),
 \qquad y=V+u.
\]
Here \(\deg U=6\), \(\deg V=5\), and the supplied gap \(u\) has degree
one.  The gap residual
\[
 (V+u)^2-Uu(2V+u)-1
\]
has degree 12, giving degree 24 after squaring.  We examined whether
canonical congruences permit a cheaper replacement of the gap.

For the canonical odd Pell index \(R\equiv3\pmod4\),
\(y=\psi_I(R)\equiv-1\pmod{I^2}\).  Thus
\(t=(y+1)/U\) is a positive integer.  The exact identity
\begin{align*}
 &U V^2-(U-1)(Ut-1)^2-1\\
 &\hspace{12mm}=
 U\bigl[V^2-U(U-1)t^2+2(U-1)t-1\bigr]
\end{align*}
gives a valid alternative residual when completeness uses these canonical
witnesses.  Its degree is 14, because of \(U(U-1)t^2\), so its square has
degree 28.  This useful divisibility does not lower the degree of the
current gap construction.  Supplying \(Uu\) instead leads to denominator
clearing and likewise did not improve the measured degree.

There is a second canonical congruence.  Write
\(A=a+2\), so \(\Delta=A^2-1\), and let
\(D=\chi_A(R)\), \(c=\psi_A(R)\).  The odd-subsequence Pell
identities, together with
\(I^2\equiv1-A^2\pmod{f^2}\), give
\[
 Ay\equiv-D\pmod{f^2}
 \qquad(R\equiv3\pmod4).
\]
For example, at \(R=3\) one has
\(\psi_I(3)=4I^2-1\), and multiplication of its reduction modulo
\(f^2\) by \(A\) gives \(-\chi_A(3)\).  The general identity follows
from the same odd-subsequence recurrence.  Notice that reducing the norm
alone gives only \((Ay)^2\equiv D^2\pmod{f^2}\); the negative branch
cannot be selected for arbitrary supplied solutions from that square
congruence alone.

The canonical congruence suggests
\(y=(f^2t-D)/A\).  Write \(V=f^2W-c\) and center the new quotient by
\(t=AW+v\).  Formally retaining \(D\) as a register, after clearing the
denominator \(A^2\) and using \(D^2-\Delta c^2=1\), the norm numerator
is divisible by \(f^2\).  The resulting expression still contains
\(-2A\Delta Wv f^4\).  This term has degree 14 under
\(\deg A=2\), \(\deg\Delta=4\), \(\deg W=3\) and
\(\deg v=\deg f=1\).  Moreover, source 58 reconstructs \(D\) rather
than supplying it as a separate unknown, so that reconstruction would also
have to be accounted for in a literal polynomial.  This calculation does
not provide the hoped-for degree-22 equation.  It proves no broader
obstruction to using this congruence in a different construction.

\subsection{Scope of the reported advance}

The demonstrated improvements in this report concern the explicit
constructions and comparisons proved in its preceding sections.  The
alternative proof substrates above remain research proposals.  They are
not counted as delivered equations, and the calculations that failed to
improve degree are retained to make the search boundary clear.  In
particular, we make neither a global simplicity-minimum claim nor a claim
that a construction independent of K215 or Jones is impossible.

Throughout, a proof of
\(\mathrm{ZFC}\vdash[\text{no zeros}\leftrightarrow
\operatorname{Con}(\mathrm{ZFC})]\)
is an internal equivalence theorem.  It is not a proof in ZFC of either
side.  Finite coefficient, identity and degree audits establish their
stated finite properties; they do not establish the consistency assertion
by failing to find a zero.


\section{Reproduction, provenance, and integration}
\label{sec:repro}

From the archive root run
\begin{quote}
\code{python3 checks/verify_all.py}
\end{quote}
to rebuild native coefficients and the nine-variable specializations
and replay the independent source, circuit and supplement audits.
Mathematical code uses only the Python standard library.
No layout search is needed: the supplied finite certificates are
the reproduction inputs. Exact polynomial checks are distinguished
from modular diagnostic evaluations.

\begin{center}
\small
\begin{tabular}{@{}p{0.46\linewidth}p{0.49\linewidth}@{}}
\toprule
Path & Purpose\\
\midrule
\path{research_nine/compact_candidate/} &
Recommended equation, builder, circuit audits and degree certificates.\\
\path{research_nine/candidate/} &
Separate original-tuple specialization.\\
\path{research_alternatives/check_positive_gap.py} &
Independent reconstruction from the authenticated source DAG and
all target interactions.\\
\path{research_low_degree/audit_gap_source.py} &
Exact \(89\to87\to86\) source identity audit.\\
\path{research_low_degree/supplement15/} &
Degree-24/89-operation companion and comparator.\\
\path{upstream44/native/} &
Complete compact tuple, native proof, pinned sources and licenses.\\
\path{upstream31/} &
Source-31 data, inherited proofs and original coefficient tuple.\\
\path{upstream-status/} &
Pinned proof-status documents and external K215 endpoint statement.\\
\path{MANIFEST.sha256} &
Hashes of distributed files, excluding the manifest itself.\\
\bottomrule
\end{tabular}
\end{center}

Imported proofs retain historical wording and relative links.
Their former frontier claims or coefficient metadata do not override
this report's compact binding. In particular, the retained
original-tuple parameters \(a_\ast=2^{2134939}\),
\(c_\ast=2^{17079501}\) appear in the general source proof and
comparison. The recommended compact parameters are
\(2^{1182967}\), \(2^{9463725}\).
The archive does not replace any certified Lean endpoint or claim
a repository Lean build.

\subsection{What would constitute a stronger next result}

Three concrete directions remain. Further positive coordinate
redesign could reduce the active pair constraints, or remove a
witness with a proved positive inverse. A direct compiler for sealed
ZFC proof codes could address the outstanding internalization
problem without the native machine layer, but must emit literal
coefficients and an internal correctness argument. A different
ratio or finalization mechanism could reduce the large numerical
coefficients in \(816\Kcode+1224W+5044\), which the present
packing improvements leave unchanged.

None of these is asserted as completed. Failing to find a zero,
checking finite coefficients, or proving a standard-natural
equivalence does not establish the internal consistency statement.
The remaining formal target must account for the internal natural
numbers of arbitrary ZFC models, including nonstandard witnesses.

\appendix
\section{Complete quadratic source}
\label{app:source}
In the following table, \(\mathcal B=B_{\mathrm{aux}}\) in the
coordinate list \eqref{eq:source-order}. The generated table is
checked against every formal coefficient in the literal source.
Its indices \(0,\ldots,30\) correspond in order to the one-based
weight list \(w_1,\ldots,w_{31}\) in the main text.
% Generated by write_source_table.py after exact formal-port comparison.
% Requires amsmath, longtable and needspace. This is an input fragment.
\subsection{The exact positive quadratic source}
All 31 variables in this table range over the positive integers.  The four symbols
$\beta,K_0,M_C,M_F$ are the coherent compact coefficient tuple; $M_F$ denotes
the port named \texttt{MF\_source} in its native binding.  Write
\[q=\beta J+1,\qquad k=\eta+\zeta.\]
The source system is $\rho_r=0$ for $1\leq r\leq22$.  The last column gives
the target exponent $\lambda_r$ assigned to that residual in the frozen packing layout.
These are local source symbols: for example, $a,c,D,N$ in this table do not denote
the outer encoding constants or the Jones-tail registers.
\begingroup
\renewcommand{\arraystretch}{1.22}
\begin{longtable}{@{}r p{0.73\linewidth} r@{}}
\hline
$r$ & $\rho_r$ & $\lambda_r$\\
\hline
\endfirsthead
\hline
$r$ & $\rho_r$ & $\lambda_r$\\
\hline
\endhead
1 & \(S-q^2\) & 906\\
2 & \(N-Sq\) & 27\\
3 & \(X-wq\) & 1499\\
4 & \(a-E-sN\) & 1255\\
5 & \(E-(a-E)X\) & 1733\\
6 & \(c-\eta-k(a-E)\) & 2016\\
7 & \(p-E(c-\eta)\) & 561\\
8 & \(t^2+2pt-kp-1\) & 1472\\
9 & \(D-X-ca-(4a+3)\gamma\) & 270\\
10 & \(\Delta-(a+1)(a+3)\) & 893\\
11 & \(C_2-c^2\) & 887\\
12 & \(D^2-C_2\Delta-1\) & 616\\
13 & \(k+V-hE-jc-1\) & 1139\\
14 & \(\begin{aligned}[t] &jc-V+\alpha(1-S)+F(N-S-q+1)\\ &\quad +(q-1)(q-N)-J(M_C+M_Fq)\end{aligned}\) & 1822\\
15 & \(\begin{aligned}[t] &K_0(q-F-\alpha)+X-F-w(F+\alpha)\\ &\quad +(q-1)(1-z)\end{aligned}\) & 966\\
16 & \(F_2-f^2\) & 1279\\
17 & \(c+V-fo\) & 686\\
18 & \(g+V^2-y^2\) & 72\\
19 & \(I-iC_2\) & 448\\
20 & \(\mathcal B-I^2\) & 320\\
21 & \(\mathcal B+(1-F_2)\Delta\) & 1190\\
22 & \(y^2-g\mathcal B-1\) & 1545\\
\hline
\end{longtable}
\endgroup
The variable $z$ in row 15 is the original transport variable.  The shifted
positive transport coordinate used by the native source is $z_{\mathrm s}=z-1$;
the positive inverse requires the source-kernel bound $z\geq2$.  The coefficient
$M_F$ retains its native binding in this change of variable.
The new coordinates are $p=P-(\eta+\zeta)$, $t=\tau-p$, and $S=L_{\rm old}+1$.
Rows 14 and 15 include the invertible row operations
$\rho_{14}=r_{14}+r_2$ and $\rho_{15}=r_{15}+r_3$.
\Needspace{16\baselineskip}
\paragraph{Variable indices and weights.}
The monomial records use the following zero-based indices and code weights $w_i$.
The unique zero weight is attached to $y$.
\begin{center}
\begin{tabular}{@{}r l r@{\qquad}r l r@{\qquad}r l r@{}}
\hline
$i$ & variable & $w_i$ & $i$ & variable & $w_i$ & $i$ & variable & $w_i$\\
\hline
0 & $J$ & 6 & 11 & $t$ & 87 & 22 & $D$ & 1\\
1 & $F$ & 10 & 12 & $\eta$ & 96 & 23 & $C_2$ & 67\\
2 & $\alpha$ & 37 & 13 & $\zeta$ & 80 & 24 & $S$ & 21\\
3 & $z$ & 46 & 14 & $y$ & 0 & 25 & $F_2$ & 59\\
4 & $f$ & 4 & 15 & $\gamma$ & 41 & 26 & $V$ & 18\\
5 & $h$ & 83 & 16 & $N$ & 3 & 27 & $g$ & 19\\
6 & $i$ & 60 & 17 & $E$ & 89 & 28 & $I$ & 61\\
7 & $o$ & 50 & 18 & $p$ & 95 & 29 & $\mathcal B$ & 14\\
8 & $j$ & 65 & 19 & $c$ & 100 & 30 & $\Delta$ & 64\\
9 & $s$ & 45 & 20 & $a$ & 97 &  &  & \\
10 & $w$ & 20 & 21 & $X$ & 84 &  &  & \\
\hline
\end{tabular}
\end{center}
\paragraph{Weighted-record convention.}
For a monomial $x^\alpha$ with $|\alpha|\leq2$, the stored coefficient is
$(2-|\alpha|)!\,\alpha!$ times its coefficient in $\rho_r$.  Thus constants,
linear terms, and diagonal quadratic terms are doubled; mixed quadratic
terms are unchanged.  The source has 86 weighted records.  The script
\texttt{write\_source\_table.py} verifies every displayed residual exactly over
$\mathbb Z[\beta,K_0,M_C,M_F]$ before writing this table.


\section{Payload identities and source authentication}

The compact raw payload SHA-256 values are:
\begin{quote}\small
\(\beta\):\\
\path{0664a868cbfef9869ea05c40a1c445563f685ff857092594c1c1223bc4b21f9d}\\
\(K_0\):\\
\path{c6d523bd32af8dd0345db843cdadf6932bfcad186912a6d112314b60657c554e}\\
\(M_C\):\\
\path{6007cabd5f5228618e8fb613f386d8fc3692a24520d6e4cd72af0262aae1ade4}\\
\(M_F\):\\
\path{48333b0d30a4edeb8579f44286cb59806c5622cb42b367f5b5f9ee827711c737}.
\end{quote}
The new 86-record source has SHA-256
\begin{quote}\small
\path{98a3be433694b5c5b5e99020c503cbb2e2a5adf05757b6bbc2357832472700ab}.
\end{quote}
The current authenticated 60-operation source has SHA-256
\begin{quote}\small
\path{ac2a67782d344d1e74fd93bc3f09fae0097eaed1829fd8194ae235ae84763e48},
\end{quote}
and Git blob \code{518b2e3f03c354ce2ecc74f8f7a75e526cd911c7}.
The historical package-59 parent hash is not claimed for the new
record format. Independent exact reconstruction from the current
source removes any reliance on an unavailable historical parent file.

\section{Reference map}

The following sources fix the scope of inherited claims. Archival
reference copies and licenses are included; full transitive Lean
dependencies remain in ProveIt.
\begin{enumerate}
\item ProveIt, reviewed snapshot \basecommit,
\emph{Computability/DiophantineConsistency}, especially
\code{OBJECT_ZFC.md}, \code{PROPOSAL_REVIEW.md} and the research collection:
\url{https://github.com/VladimirReshetnikov/ProveIt/tree/beb0192cb8dc36f78c369c080f2adfbc1887a486/Computability/DiophantineConsistency}.
\item Package 59, \emph{Free-coordinate composition with the
degree-separated source}, \path{nine_free_adjunct/PROOF.md}, and
its retained FreeCoordinate proof and review: numerical tail
and comparison metrics.
\item Package 31, \code{quartic31/QUARTIC60_PROOF.md},
\code{quartic31_system60.json}, and ordinary compiler/Pell proofs:
quadratic source and positive transport shift.
\item Package 44, \code{native/PROOF.md},
\code{COEFFICIENT_BINDING.json}, \code{SOURCE_PINS.json},
\code{build_coefficients.py} and \code{audit_coefficients.py}:
whole compact tuple and native semantic adapter.
\item Package 58, report and \code{code/build.py}: degree-24 parent.
\item \code{BusyBeaver.ZFC.KEndpoint215},
\code{blank_halts_iff_zfc_refutation}: external literal-machine
endpoint, distinct from the internal target in
\eqref{eq:internal-target}.
\item J. P. Jones, \emph{Universal Diophantine equation},
Journal of Symbolic Logic 47 (1982), 549--571. The exact numerical
ratio interfaces here are the repository's literal versions.
\item J. Bayer and M. David, \emph{A Formal Proof of Complexity Bounds
on Diophantine Equations}, ITP 2025, LIPIcs 352, 3:1--3:18;
\url{https://arxiv.org/abs/2505.16963v2}. Associated AFP development:
\url{https://isa-afp.org/entries/Diophantine_Universal_Pairs.html}.
\end{enumerate}

\end{document}
